\begin{choice}[topic={全称命题的否定}]
若命题 ``$\forall x \in \mathbb{R}, 1-x^2 \le m$'' 是假命题，则实数 $m$ 的取值范围是
\begin{tasks}(2)
	\task $(-\infty, 1)$
	\task $(-\infty, 1]$
	\task $(1, +\infty)$
	\task $[1, +\infty)$
\end{tasks}
\end{choice}
\begin{choice-sol}
\textbf{答案 A}

一个全称量词命题为假，其逻辑等价于它的否定形式（一个存在量词命题）为真.
设原命题为 $P: \forall x \in \mathbb{R}, 1-x^2 \le m$。

由于命题 $P$ 为假，其否定命题 $\neg P$ 必为真.
\[ \neg P: \exists x \in \mathbb{R}, \neg(1-x^2 \le m) \implies \exists x \in \mathbb{R}, 1-x^2 > m \]
该存在量词命题为真，意味着 $m$ 必须小于函数 $f(x) = 1-x^2$ 的最大值.
函数 $f(x)=1-x^2$ 是一个开口向下的二次函数，其顶点为最大值点.
\[ f(x)_{\max} = f(0) = 1 \]
因此，实数 $m$ 的取值范围必须是 $m < 1$，
即 $m \in (-\infty, 1)$。
\end{choice-sol}
